\documentclass[../../main.tex]{subfiles}

\begin{document}

\chapter{[A] Functional Equations}

Personally I think one of the most difficult sections in olympiad
algebra is functional equations. Unlike other sections that involves
application of background knowledge, functional equation problems
tend to prioritize ingenuity and creativity. This could be a
good news if you are cramming for an exam, but I think it is
important to solve many problems to find different patterns in
problems involving functional equations.

\section{Foundations}

Before we get into what a functional equation is, let's briefly
review the terms related to functions as they appear quite a lot.

\begin{minipage}{0.48\textwidth}
\begin{center}
\textbf{Injective Function}
\begin{align*}
    &\text{If } f(x_1) = f(x_2), \\
    &\text{then } x_1 = x_2
\end{align*}
\end{center}
\end{minipage}
\hfill
\begin{minipage}{0.48\textwidth}
\begin{center}
\begin{tikzpicture}[
    >=latex, very thick,
    inout/.style={
        draw, fill=BrickRed!10, ellipse,
        minimum width=1.3cm, minimum height=3cm
    },
    labelnode/.style={font=\large, black},
    scale=0.60
]

    \node[inout] (L) at (0,0) {};
    \node[labelnode, above=2mm of L.north] {X};

    \foreach \i/\xs in {1/$x_1$,2/$x_2$,3/$x_3$}{
        \node at ($(L) + (0cm,1.6cm - \i*0.8cm)$) {\xs};
    }

    \node[inout] (R) at (4,0) {};
    \node[labelnode, above=2mm of R.north] {Y};

    \foreach \i/\ys in {1/$y_1$,2/$y_2$,3/$y_3$,4/$y_4$}{
        \node at ($(R) + (0cm,1.8cm - \i*0.8cm)$) {\ys};
    }

    \draw[->, Sepia]
        ($(L)+(0.5cm,0.8cm)$) to ($(R)+(-0.5cm,1.0cm)$);
    \draw[->, Sepia]
        ($(L)+(0.5cm,0.0cm)$) to ($(R)+(-0.5cm,-1.4cm)$);
    \draw[->, Sepia]
        ($(L)+(0.5cm,-0.8cm)$) to ($(R)+(-0.5cm,-0.6cm)$);
\end{tikzpicture}
\end{center}
\end{minipage}

\begin{minipage}{0.48\textwidth}
\begin{center}
\textbf{Surjective Function}
\begin{align*}
    &\forall\, y_i \in Y,\ \exists\, x_j \in X \\
    &\quad s.t.\ f(x_j) = y_i
\end{align*}
\end{center}
\end{minipage}
\hfill
\begin{minipage}{0.48\textwidth}
\begin{center}
\begin{tikzpicture}[
    >=latex, very thick,
    inout/.style={
        draw, fill=MidnightBlue!10, ellipse,
        minimum width=1.3cm, minimum height=3cm
    },
    labelnode/.style={font=\large, black},
    scale=0.6
]

    \node[inout] (L) at (0,0) {};
    \node[labelnode, above=2mm of L.north] {X};

    \foreach \i/\xs in {1/$x_1$,2/$x_2$,3/$x_3$,4/$x_4$}{
        \node at ($(L) + (0cm,1.8cm - \i*0.8cm)$) {\xs};
    }

    \node[inout] (R) at (4,0) {};
    \node[labelnode, above=2mm of R.north] {Y};

    \foreach \i/\ys in {1/$y_1$,2/$y_2$,3/$y_3$}{
        \node at ($(R) + (0cm,1.6cm - \i*0.8cm)$) {\ys};
    }

    \draw[->, Sepia] 
        ($(L)+(0.5cm,1.0cm)$) to ($(R)+(-0.5cm,0.8cm)$);
    \draw[->, Sepia] 
        ($(L)+(0.5cm,0.2cm)$) to ($(R)+(-0.5cm,-0.8cm)$);
    \draw[->, Sepia] 
        ($(L)+(0.5cm,-0.6cm)$) to ($(R)+(-0.5cm,0.0cm)$);
    \draw[->, Sepia]
        ($(L)+(0.5cm,-1.4cm)$) to ($(R)+(-0.5cm,-0.8cm)$);
\end{tikzpicture}
\end{center}
\end{minipage}

\begin{minipage}{0.48\textwidth}
\begin{center}
\textbf{Bijective Function}
\begin{align*}
    &\forall\, y_i \in Y,\ \exists!\, x_j \in X \\
    &\quad s.t.\ f(x_j) = y_i
\end{align*}
\end{center}
\end{minipage}
\hfill
\begin{minipage}{0.48\textwidth}
\begin{center}
\begin{tikzpicture}[
    >=latex, very thick,
    inout/.style={
        draw, fill=RoyalPurple!10, ellipse,
        minimum width=1.3cm, minimum height=3cm
    },
    labelnode/.style={font=\large, black},
    scale=0.6
]

    \node[inout] (L) at (0,0) {};
    \node[labelnode, above=2mm of L.north] {X};
  
    \foreach \i/\xs in {1/$x_1$,2/$x_2$,3/$x_3$,4/$x_4$}{
        \node at ($(L) + (0cm,1.8cm - \i*0.8cm)$) {\xs};
    }

    \node[inout] (R) at (4,0) {};
    \node[labelnode, above=2mm of R.north] {Y};

    \foreach \i/\ys in {1/$y_1$,2/$y_2$,3/$y_3$,4/$y_4$}{
        \node at ($(R) + (0cm,1.8cm - \i*0.8cm)$) {\ys};
    }
  
    \draw[->, Sepia]
        ($(L)+(0.5cm,1.0cm)$) to ($(R)+(-0.5cm,1.0cm)$);
    \draw[->, Sepia]
        ($(L)+(0.5cm,0.2cm)$) to ($(R)+(-0.5cm,-1.4cm)$);
    \draw[->, Sepia]
        ($(L)+(0.5cm,-0.6cm)$) to ($(R)+(-0.5cm,0.2cm)$);
    \draw[->, Sepia]
        ($(L)+(0.5cm,-1.4cm)$) to ($(R)+(-0.5cm,-0.6cm)$);
\end{tikzpicture}
\end{center}
\end{minipage}

While we will be using the terms above, it is worth noting other
famous terms describing the same functions.

\begin{important}
Some people use the term \textbf{one-to-one} for injective,
\textbf{onto} for surjective, and \textbf{one-to-one onto} for
bijective functions.
\end{important}

As mentioned, we don't really need other strong backgrounds for
functional equations, though there are plenty of techniques to study.
However, the terms injective, surjective, and bijective really cannot
be stressed enough.

\begin{important}
A function $f$ has an inverse if and only if it is bijective.
\end{important}

This seemingly obvious statement can come very handy. To show it,
you can consider the diagrams above. If the function is only
injective, then we have inputs with no corresponding output.
If $f$ is surjective only, then there are inputs with two outputs,
which is not a function anymore. Another very interesting
property to note is strictly increasing and decreasing functions.

\begin{important}
A strictly increasing or decreasing function is always injective.
\end{important}

First, note that a strictly increasing function $f$ have the property
where if $x < y$ then $f(x) < f(y)$. This satisfies the definition
of injective function because $f$ is \textit{strictly} increasing.
The case for strictly decreasing is proven in analogous way.
With these in mind, let's get started with functional equations!

A functional equation as the name suggest is simply an equation
where the variable is a function. To see what it means, let's
solve practice problems. Really, there is nothing more to it
and we now just need to keep practicing pattern recognition and
thinking creatively. However before starting, there is an important
note to keep in mind.

\begin{important}
The phrase \textbf{if and only if} cannot be stressed enough!
If you think you've found all the solutions, chances are you
probably haven't. Many points are taken off because of the lack
of the proof for the \textbf{if and only if} condition, and it
is always nice to suspect your assumption that you have found all
the solutions.
\end{important}

For our first problem, let's solve a classic problem that involves
injection, surjection, and bijection.

\begin{problem}{Classic Problem \easy}{}
Find all functions $f$ such that $f(f(x) + f(y)) = x + y$ for all
$x, y \in \mathbb{Z}$ \cite{MATH01-5}.
\end{problem}
\begin{solution}
Unlike the traditional approach of substituting different values
to $x$ and $y$, let's start by substituting $y = 1$. Substituting,
$f(f(x) + f(1)) = x + 1 = f(f(x + 1) + f(0))$. To find injectivity,
surjectivity, or bijectivity of the function $f$, assume
$f(a) = f(b)$. Then, $f(f(a) + f(y)) = f(f(b) + f(y))$ and by
definition, $f(a) = f(b)$ leads to $a = b$. In other words,
$f$ is injective.

Returning to our previous equation $f(f(x) + f(1)) =
f(f(x + 1) + f(0))$, we could see that $f(x) + f(1) = f(x + 1) + f(0)$
since $f$ is injective. Rearranging, $f(x + 1) - f(x) = f(1) - f(0)$
and $f(x)$ is arithmetic sequence. Therefore let $f(x) = ax + b$.
Substituting to our given equation, $f(ax + b + ay + b) =
a (ax + b + ay + b) + b = x + y$. Solving the equation,
\begin{align*}
    a (ax + b + ay + b) + b &= x + y \\
    \left( a^2 - 1 \right) x + \left( a^2 - 1 \right) y
        + 2ab + b &= 0
\end{align*}
holds for all $x, y \in \mathbb{R}$. Thus $a = \pm 1$ and $b = 0$.
Therefore, $f(x) = \pm x$ are all solutions.
\end{solution}

Continuing from this problem, we have a harder problem that requires
smart substitutions.

\begin{problem}{2008 IMO Problem 4 \hard}{2008 IMO Problem 4}
Find all functions $f: (0, \infty) \mapsto (0, \infty)$ that satisfy
the following equation for all positive reals $a, b, c, d$ where
$ab = cd$.
\[
\frac{ f(a)^2 + f(b)^2 }{f(c^2) + f(d^2)}
= \frac{a^2 + b^2}{c^2 + d^2}
\]
\end{problem}
\begin{solution}
First, let's start with what we can observe from the problem.
The first thing that we can see is that we have squares everywhere.
The second thing to note is that substituting $a = b = c = d = 1$
gives us $f(x) = 1$. We can also notice that if we substitute
$(n, n, n, n)$ we get $f(n)^2 = f(n^2)$.

With the observations in mind, what we can try is getting rid of
the square by substituting $(\sqrt{x}, \sqrt{x}, 1, x)$. This
will lead us to following equations.
\begin{align*}
    \frac{ f(\sqrt{x})^2 + f(\sqrt{x})^2 }{ 1 + f(x^2) }
    &= \frac{x + x}{1 + x^2} \\
    \frac{ 2f(x) }{ 1 + f(x)^2 }
    &= \frac{2x}{1 + x^2} \\
    2f(x) (1 + x^2)
    &= 2x(1 + f(x)^2)
\end{align*}
Continuing, we can rearrange the equations to solve for $f$.
\begin{align*}
    xf(x)^2 - (1 + x^2)f(x) + x &= 0 \\
    (xf(x) - 1) (f(x) - x) &= 0 \\
    \therefore
    f(x) = x \text{ or } f(x) &= \frac{1}{x}
\end{align*}
Now we can see that if a function $f$ satisfy the equation given
by the problem, then $f(x) = x$ or $f(x) = \frac{1}{x}$. Moreover,
if $f(x) = x)$ or $f(x) = \frac{1}{x}$, then $f$ satisfies the
original equation as we can see from direct substitution.

If we end our solution here thinking that we have solved the problem,
we will lose points for not showing that $x$ and $\frac{1}{x}$
are the only solutions. Although we showed the if and only if
condition, we have not showed that $f$ is not a piece-wise function.

\begin{lemma*}
$f(x) = x$ and $f(x) = \frac{1}{x}$ are the only solutions to the
equation.
\end{lemma*}
\begin{proofInBox}
For the sake of contradiction, let $f(\alpha) \neq \alpha$ and
$f(\beta) \neq \frac{1}{\beta}$ for $\alpha, \beta \neq 1$.
By our previous conclusions, we can see that it implies
$f(\alpha) = \frac{1}{\alpha}$ and $f(\beta) = \beta$.

Substitute $(\sqrt{\alpha}, \sqrt{\beta}, 1, \sqrt{\alpha\beta})$
to eliminate the squares and use both $\alpha$ and $\beta$.
\begin{align*}
    \frac{ f(\sqrt{\alpha})^2 + f(\sqrt{\beta})^2 }
        { f(1) + f(\alpha\beta) }
    &= \frac{\alpha + \beta}{1 + \alpha\beta} \\
    \frac{ \frac{1}{\alpha} + \beta }{ 1 + f(\alpha\beta) }
    &= \frac{\alpha + \beta}{1 + \alpha\beta}
\end{align*}
Notice that $f(\alpha\beta)$ is either $\alpha\beta$ or
$\frac{1}{\alpha\beta}$. Therefore, let $f(\alpha\beta) =
\alpha\beta$. Solving the equation, the following equations are
obtained.
\begin{align*}
    \frac{ \frac{1}{\alpha} + \beta }{ 1 + f(\alpha\beta) }
    &= \frac{\alpha + \beta}{1 + \alpha\beta} \\
    \frac{ \frac{1}{\alpha} + \beta }{ 1 + \alpha\beta }
    &= \frac{\alpha + \beta}{1 + \alpha\beta} \\
    \frac{1}{\alpha} &= \alpha
\end{align*}
This contradicts with the fact that $\alpha \neq 1$. Continuing,
let $f(\alpha\beta) = \frac{1}{\alpha\beta}$.
\begin{align*}
    \frac{ \frac{1}{\alpha} + \beta }{ 1 + f(\alpha\beta) }
    \displaybreak
    &= \frac{\alpha + \beta}{1 + \alpha\beta} \\
    \frac{ \frac{1}{\alpha} + \beta }{ 1 + \frac{1}{\alpha\beta} }
    &= \frac{\alpha + \beta}{1 + \alpha\beta} \\
    \frac{\beta + \alpha\beta^2}{\alpha\beta + 1}
    &= \frac{\alpha + \beta}{1 + \alpha\beta} \\
    \alpha\beta^2 &= \alpha \\
    \beta^2 &= 1
\end{align*}
This also contradicts with $\beta \neq 1$. Therefore, there exists
no $\alpha, \beta \neq 1$ such that $f(\alpha) \neq \alpha$ and
$f(\beta) \neq \frac{1}{\beta}$.
\end{proofInBox}

Hence, $f(x) = x$ and $f(x) = \frac{1}{x}$ are the only solutions.
\end{solution}

As you probably have noticed from the solution, while vast amount
of knowledge in different functions may come handy, ingenuity
and clever substitution would do most of the work when solving
problems. Here is another problem on using clever substitutions.

\begin{problem}
{2020 Baltic Way Problem 4 \easy}
{2020 Baltic Way Problem 4}
Find all functions $f: \mathbb{R} \rightarrow \mathbb{R}$ that
satisfy the following for all $x, y \in \mathbb{R}$.
\[
f(f(x) + x + y) = f(x + y) + yf(y)
\]
\end{problem}
(\href{https://youtu.be/STy74sON1Ho}{Video Solution})
\begin{solution}
First, let's start by doing typical substitutions. Substituting
$x = y = 0$ gives $f(f(0)) = f(0)$ and $x = -y$ gives
$f(f(x)) = f(0) - xf(-x)$. Also, trying $y = 0$ gives $f(f(x) + x)
= f(x)$. Looking at the last observation, we could assume that
$f(x) = 0$. Moreover our first observation also suggests that
$f(0) = 0$. With the observations, let's first investigate the
value of $f(0)$.

Consider the following equation by the last observation with
$y = f(x) + x$.
\begin{align*}
    f(x)
    &= f(y)
    = f(f(y) + y)
    = f(f(x) + f(x) + x) \\
    &= f(x + f(x)) + f(x) f(f(x))
    = f(x) + f(x) f(f(x))
\end{align*}
Therefore, $f(x) f(f(x)) = 0$ and $f(0) f(f(0)) = 0$. Note that
because $f(f(0)) = f(0)$, $f(0)^2 = 0$ and $f(0) = 0$ holds.
Indeed $f(0) = 0$!

Continuing, substitute $x = 0$ to the given equation.
\[
f(f(0) + 0 + y)
= f(y)
= f(0 + y) + y f(y)
\]
In other words, $y f(y) = 0$ holds for all $y \in \mathbb{R}$.
Specially, $f(y) = 0$ for all $y \neq 0$ and we have found that
$f(0) = 0$. Therefore, $f(x) = 0$ for all $x \in \mathbb{R}$ is
the only solution.
\end{solution}

In this problem, using the observation carefully and letting
$y = f(x) + x$ was the key part of the solution. There are
many topics to discuss even if it requires less prerequisite
knowledge to actually solve the problem. With the intuition
discussed in this note, let's discuss few important topics and
patterns that are widely used.

\end{document}


